% LSPV-1 REPORT::PAIR2COVER1::2026-10-01T15:07:00-07:00
% !TEX program = xelatex
\documentclass[11pt]{article}
\usepackage[margin=0.88in]{geometry}
\usepackage{amsmath,amssymb}
\begin{document}
\section*{LSPV-1: pair-cover criterion for global Hadamard quadratic rank}
\textbf{STATUS: OPEN.} P versus NP remains open. This report proves an exact structural criterion for the declared SU05-sanctioned Hadamard realization. It proves no unrestricted SAT lower bound and does not prove that the active vMS legal-state family satisfies the criterion.

\paragraph{Recovery preflight.}
The newest verified LSPV-1 artifact visible at preflight remains \texttt{2026-10-01\_0915\_LSPV-1\_HADCOEF1.tex}. Later failed cycles are described in run context, but their exact failed LaTeX bodies are not available as repository artifacts. I therefore did not reconstruct or duplicate them. GitHub is readable and peers have later verified writes, so absent the original connector error the immediate historical cause remains \textbf{UNKNOWN}.

\paragraph{CORPUS REVIEWED.}
I surveyed the current repository tree and search-visible history; read root \texttt{main.tex}, \texttt{AGENTS.md}, \texttt{CODEX.md}, the repo-scoped pnp-research skill, TASK-001 and the current task index; inspected newest/relevant MAIN material and materially relevant LSPV reports, including LSPV-1 SCHUR1/LEAF2/MASK2/MASKBASIS1/ENCSTOP1/SU05CODE1/HADCOEF1, LSPV-2 HADSQ1 and POLAR1, LSPV-3 DEPSWAP1, and LSPV-4 PCLOSE1-R1; surveyed the random-challenge HEC/CRS-PPC/CCSPC branch, governance/corpus-reconciliation rules, experiment/result paths referenced by root MAIN, and repository search hits for CLAUDE/GROK and prior product-span/common-kernel formulations. Root LSPV-1 still foregrounds the older leaf-full-rank target and HADCOEF1. Verified POLAR1 supplies the exact common-isotropic-slice formulation used here.

\paragraph{Primary-literature preflight and two obstacles.}
The first obstacle is exact typing of the legal reconstruction-compatible SU/vMS state family and its Boolean label maps. Shaltiel--Umans (2005) supplies the algebraic extractor and permits the Hadamard conversion in the relevant computational conversion, while van Melkebeek--Mocelin Sdroievski (2025) supplies the PCPP/RMV targeted-HSG architecture but no Boolean Schur-square/common-kernel theorem. The second obstacle is global transversality: small fixed-state Hadamard slices can have large kernels whose intersection is nevertheless zero. Componentwise products of linear codes are exactly the relevant coding-theory object (Randriambololona), but no reviewed coding theorem supplies the missing SU/vMS legal-state incidence statement.

\paragraph{Setup.}
For the declared Hadamard conversion, let each legal source state \(v\) determine the linear output slice
\[
S_v=\operatorname{im}X_v\le \mathbb F_2^m.
\]
Let
\[
\mathcal Q_2=\left\{q(y)=c+\ell(y)+\sum_{i<j}\beta_{ij}y_i y_j\right\},
\qquad M=\dim\mathcal Q_2=1+m+\binom m2.
\]
POLAR1 gives exactly
\[
K_2(H)=\{q\in\mathcal Q_2:q|_{S_v}=0\text{ for every legal }v\}.
\]
For \(q(0)=0\), define \(b_q(x,y)=q(x+y)+q(x)+q(y)\).

\paragraph{PROVED/EXACT: pair-cover criterion.}
Suppose there exist \(t_1,\ldots,t_m\in\mathbb F_2^m\) such that
\[
\operatorname{span}\{t_1,\ldots,t_m\}=\mathbb F_2^m,
\]
and
\[
\forall i\ \exists v: t_i\in S_v,
\qquad \forall i<j\ \exists v: t_i,t_j\in S_v.
\]
Then
\[
\boxed{K_2(H)=\{0\}},\qquad
\boxed{\rho(H)=M=1+m+\binom m2}.
\]
For \(q\in K_2(H)\), every slice contains zero, so \(q(0)=0\). Singleton coverage gives \(q(t_i)=0\). Pair coverage gives a legal linear slice containing \(t_i,t_j,t_i+t_j\), hence \(b_q(t_i,t_j)=0\). Choose a basis among the spanning \(t_i\). The polarization identity
\[
q\!\left(\sum_i\lambda_i t_i\right)=
\sum_i\lambda_i q(t_i)+
\sum_{i<j}\lambda_i\lambda_j b_q(t_i,t_j)
\]
gives \(q\equiv0\) on all \(\mathbb F_2^m\). Thus the global degree-two evaluation map has full rank. \(\square\)

\paragraph{Equivalent incidence formulation.}
Define the legal-slice co-occurrence graph on nonzero vectors appearing in at least one \(S_v\): join \(x,y\) when some legal slice contains both. The theorem requires a spanning set of \(m\) vectors inducing a clique. This is weaker than requiring one legal slice of dimension \(m\), and avoids the invalid rectangularization step of replacing same-source pairs by an arbitrary Cartesian product. Cross terms are certified only by actual legal co-occurrence.

\paragraph{PROVED/EXACT: partial pair-cover lower bound.}
Let \(T=\{t_1,\ldots,t_r\}\) be linearly independent, assume each \(t_i\) occurs in a legal slice, and let \(E\subseteq\binom{[r]}2\) be the pairs co-contained in at least one legal slice. The equations \(q(t_i)=0\) isolate the \(r\) linear coefficients in coordinates relative to \(T\), while \(b_q(t_i,t_j)=0\) for \(\{i,j\}\in E\) isolate distinct quadratic coefficients. With \(q(0)=0\), these are independent functionals on \(\mathcal Q_2\). Therefore
\[
\boxed{\rho(H)\ge 1+r+|E|}.
\]
A clique on \(r\) independent vectors recovers \(1+r+\binom r2\), but the clique may be assembled across many legal states.

\paragraph{Consequence for the current target.}
The root target that every nonconstant univariate leaf be full rank remains sufficient but is stronger than necessary. To refute a global \(O(m)\) product-span theorem it suffices to prove, for one SAME\_REALIZATION legal family,
\[
r+|E|=\omega(m).
\]
To prove maximal rank it suffices to find a spanning \(m\)-clique. Conversely, an \(O(m)\) theorem forces \(r+|E|=O(m)\) for every independent legal direction set, a concrete architecture-specific sparsity requirement.

\paragraph{Boolean source degree.}
Nothing here bounds the ANF degree of \(v\mapsto x_{i,a}(v)\). Odd-characteristic polynomial degree of the SU/RMV source map still does not imply low Boolean degree after label conversion. The source-degree arrow remains \textbf{OPEN}; field-degree and locality shortcuts remain retired.

\paragraph{PQ/cross-domain classification.}
\textbf{EXACT OBJECT/THEOREM TRANSFER:} Schur/componentwise products, quadratic-relation kernels, Boolean polarization, fixed-code puncturing, and legal-slice co-occurrence are exact for the declared realization.
\textbf{ANALOGY/CANDIDATE:} shortening filtrations, minimum relation rank, Hilbert functions, syzygies/graded Betti numbers, matroid/incidence expansion, and LWE/special-soundness extraction architecture require new construction-specific bridges.
\textbf{FAILED TRANSFER:} cryptographic hardness/random-oracle assumptions, VLSI area-time, control barriers, distributed cuts, topology, and proof-system-specific lower bounds do not yield unrestricted random-access SAT runtime.

\paragraph{Hostile controls and barriers.}
SAME\_REALIZATION fixes the legal RMV restriction, state family, encoder, mask family, and consumer. A semantic oracle may supply output semantics and an algebraic oracle may expose all \(S_v\) and their relation space, so this theorem is structural, relativizing, and algebraic. DEAD/COPY/DELAY preserve incidence/rank while changing traces; INTERP/RECOMP or nonlinear recoding can change it and must be retyped. Padding is not runtime. Affine/Tseitin with Gaussian elimination, Horn, 2-SAT, bounded treewidth, independent products, and Generic Avoid remain hostile controls: dense pair-cover or full Schur rank is not itself computational hardness.

\paragraph{DUPLICATION / CONTRADICTION AUDIT.}
HADCOEF1 already identifies \(K_2(H)=\ker\Psi\). HADSQ1 gives exact one-slice rank. POLAR1 gives common-isotropic-slice equations. This report does not relabel those results. The new theorem is the cross-state pair-cover criterion and partial incidence lower bound. It resolves the earlier CROSSSPAN-style temptation without illicit rectangularization. No contradiction with the small-slice cap exists because many low-dimensional slices can collectively certify all pairs.

\paragraph{WHAT CHANGED RELATIVE TO ROOT MAIN.}
Root MAIN's LSPV-1 region is stale in two respects: it still makes the all-parameter leaf-full-rank statement the current target, and it does not incorporate the verified polar/common-slice formulation. The weaker and more precise target is legal pair incidence across the SAME\_REALIZATION family. Full rank of every leaf is no longer required.

\paragraph{ADOPTED / REJECTED / RETIRED PRIOR COMPONENTS.}
Adopted: source-sanctioned Hadamard typing, HADCOEF1, HADSQ1, POLAR1, fixed-code Schur squares, and charged acquisition. Rejected: generic rectangularization from same-source identities to arbitrary cross products. Retired as a necessary target: one-state high rank and every-nonconstant-leaf full rank. Retained only as sufficient evidence: older exhaustive leaf experiments. Generic short-source/canonicity/locality/field-degree/list-decoding implications remain falsified/retired.

\paragraph{Strict status and single next target.}
\textbf{PROVED/EXACT:} pair-cover theorem and partial incidence lower bound for the declared Hadamard realization.
\textbf{CONDITIONAL:} application to the active reconstruction-compatible vMS family pending exact legal \(X_v\) transcription.
\textbf{OPEN:} exact legal slice family, Boolean source degree, global rank for that family, acquisition/runtime transfer, reconstruction binding, nonrelativizing/nonalgebrizing hinge, and final diagonal.

Next: extract the earliest reconstruction-compatible legal matrices \(X_v\), form bases for \(S_v=\operatorname{im}X_v\), and build the co-occurrence graph. Search first for an independent set \(T\) with large \(|E|\), not for a single high-rank slice. A spanning clique proves maximal \(\Theta(m^2)\) rank; an architecture theorem forcing \(|E|=O(m)\) for every independent \(T\) would identify the sparsity needed by an \(O(m)\) contraction.

\begin{thebibliography}{5}
\bibitem{SU05} R. Shaltiel and C. Umans, \emph{Simple Extractors for All Min-Entropies and a New Pseudorandom Generator}, J. ACM 52(2):172--216, 2005.
\bibitem{vMS25} D. van Melkebeek and N. Mocelin Sdroievski, \emph{Instance-Wise Hardness and Refutation versus Derandomization for Arthur--Merlin Protocols}, Computational Complexity 34:15, 2025.
\bibitem{RandProd} H. Randriambololona, \emph{On products and powers of linear codes under componentwise multiplication}, 2013, arXiv:1312.0022.
\end{thebibliography}
\end{document}
% END_OF_VERIFIED_REPORT::LSPV-1::PAIR2COVER1